% Supplementary G — Journal §VII (Discussion)
% Limitations and threats to validity at depth.

\documentclass[11pt,a4paper]{article}
% Shared preamble for all supplementary chapter documents.
% Used by each supplementary/conference/suppl_*.tex and
% supplementary/journal/suppl_*.tex via % Shared preamble for all supplementary chapter documents.
% Used by each supplementary/conference/suppl_*.tex and
% supplementary/journal/suppl_*.tex via % Shared preamble for all supplementary chapter documents.
% Used by each supplementary/conference/suppl_*.tex and
% supplementary/journal/suppl_*.tex via \input{../preamble.tex}.

\usepackage[utf8]{inputenc}
\usepackage[T1]{fontenc}
\usepackage[a4paper,margin=2.2cm]{geometry}
\usepackage{amsmath,amssymb,amsfonts,amsthm}
\usepackage{mathtools}
\usepackage{graphicx}
\usepackage{booktabs}
\usepackage{array}
\usepackage{multirow}
\usepackage{longtable}
\usepackage{algorithm}
\usepackage{algpseudocode}
\usepackage{xcolor}
\usepackage{url}
\usepackage{cite}
\usepackage[hidelinks]{hyperref}

\graphicspath{{../../figures/}}

\renewcommand{\thesection}{\Alph{section}.\arabic{section}}
\renewcommand{\thesubsection}{\thesection.\arabic{subsection}}

\theoremstyle{plain}
\newtheorem{theorem}{Theorem}[section]
\newtheorem{lemma}[theorem]{Lemma}
\newtheorem{proposition}[theorem]{Proposition}
\theoremstyle{definition}
\newtheorem{definition}[theorem]{Definition}
\theoremstyle{remark}
\newtheorem{remark}[theorem]{Remark}

\setcounter{secnumdepth}{3}
\setcounter{tocdepth}{2}

\newcommand{\R}{\mathbb{R}}
\newcommand{\E}{\mathbb{E}}
\newcommand{\KL}[2]{\mathrm{KL}\!\left(#1 \,\middle\|\, #2\right)}
\newcommand{\ari}{\mathrm{ARI}}
\newcommand{\nmi}{\mathrm{NMI}}
\newcommand{\softmax}{\mathrm{softmax}}
.

\usepackage[utf8]{inputenc}
\usepackage[T1]{fontenc}
\usepackage[a4paper,margin=2.2cm]{geometry}
\usepackage{amsmath,amssymb,amsfonts,amsthm}
\usepackage{mathtools}
\usepackage{graphicx}
\usepackage{booktabs}
\usepackage{array}
\usepackage{multirow}
\usepackage{longtable}
\usepackage{algorithm}
\usepackage{algpseudocode}
\usepackage{xcolor}
\usepackage{url}
\usepackage{cite}
\usepackage[hidelinks]{hyperref}

\graphicspath{{../../figures/}}

\renewcommand{\thesection}{\Alph{section}.\arabic{section}}
\renewcommand{\thesubsection}{\thesection.\arabic{subsection}}

\theoremstyle{plain}
\newtheorem{theorem}{Theorem}[section]
\newtheorem{lemma}[theorem]{Lemma}
\newtheorem{proposition}[theorem]{Proposition}
\theoremstyle{definition}
\newtheorem{definition}[theorem]{Definition}
\theoremstyle{remark}
\newtheorem{remark}[theorem]{Remark}

\setcounter{secnumdepth}{3}
\setcounter{tocdepth}{2}

\newcommand{\R}{\mathbb{R}}
\newcommand{\E}{\mathbb{E}}
\newcommand{\KL}[2]{\mathrm{KL}\!\left(#1 \,\middle\|\, #2\right)}
\newcommand{\ari}{\mathrm{ARI}}
\newcommand{\nmi}{\mathrm{NMI}}
\newcommand{\softmax}{\mathrm{softmax}}
.

\usepackage[utf8]{inputenc}
\usepackage[T1]{fontenc}
\usepackage[a4paper,margin=2.2cm]{geometry}
\usepackage{amsmath,amssymb,amsfonts,amsthm}
\usepackage{mathtools}
\usepackage{graphicx}
\usepackage{booktabs}
\usepackage{array}
\usepackage{multirow}
\usepackage{longtable}
\usepackage{algorithm}
\usepackage{algpseudocode}
\usepackage{xcolor}
\usepackage{url}
\usepackage{cite}
\usepackage[hidelinks]{hyperref}

\graphicspath{{../../figures/}}

\renewcommand{\thesection}{\Alph{section}.\arabic{section}}
\renewcommand{\thesubsection}{\thesection.\arabic{subsection}}

\theoremstyle{plain}
\newtheorem{theorem}{Theorem}[section]
\newtheorem{lemma}[theorem]{Lemma}
\newtheorem{proposition}[theorem]{Proposition}
\theoremstyle{definition}
\newtheorem{definition}[theorem]{Definition}
\theoremstyle{remark}
\newtheorem{remark}[theorem]{Remark}

\setcounter{secnumdepth}{3}
\setcounter{tocdepth}{2}

\newcommand{\R}{\mathbb{R}}
\newcommand{\E}{\mathbb{E}}
\newcommand{\KL}[2]{\mathrm{KL}\!\left(#1 \,\middle\|\, #2\right)}
\newcommand{\ari}{\mathrm{ARI}}
\newcommand{\nmi}{\mathrm{NMI}}
\newcommand{\softmax}{\mathrm{softmax}}


\title{Journal Supplementary G --- Limitations and threats to validity\\
{\large Companion to: \emph{Beyond Accuracy} (Section VII)}}
\author{Felipe Santib\'a\~nez-Leal}
\date{2026}

\begin{document}
\maketitle

\section{Scope}
This supplement enumerates the limitations and threats to validity
of the twelve-axis framework at greater depth than journal Section
VII admits. We organise them as (a)~design-time constraints baked
into the framework, (b)~empirical caveats specific to the present
benchmark set, (c)~statistical caveats around the inference, and
(d)~reporting choices that may need revision in follow-up work.

\section{Design-time constraints}

\subsection{Canonical $K$ is held fixed within an axis}
Every axis is computed at the canonical $K = \max(4, \min(12,
n_{\text{classes}}))$ except for F-4 (capacity sweep). This is by
design: changing $K$ between axes would introduce a second degree of
freedom that mixes with the axis-specific perturbation
(seed / band / preprocessing / scene). The trade-off is that the
per-axis numbers depend on a non-trivially chosen $K$ that a future
paper may want to vary; the framework admits per-axis $K$-sweeps as
a forward-compatible extension without redesign.

\subsection{Uniform tokenisation across scenes}
$Q + 1 = 13$ quantisation levels is fixed across scenes. Scenes with
narrow spectral dynamic range (Pavia U, ROSIS, 8-bit per band) and
wide dynamic range (AVIRIS, 12-bit) are tokenised at the same
resolution, which may waste resolution on the narrow-range scenes.
A future paper exploring per-scene adaptive $Q$ should re-run all
axes; the framework will automatically pick up the new $Q$ from the
canonical-fit hyperparameter slate.

\subsection{Hyperparameters $(\alpha, \eta)$ held fixed}
$\alpha = 0.45, \eta = 0.20$ across all fits. These are scikit-learn
LDA defaults that we have not tuned per scene. The companion code
repository runs a hyperparameter grid in
\texttt{lda\_hyperparam\_search/} for reference but the framework's
canonical fit deliberately fixes the slate to keep cross-axis
comparisons clean.

\subsection{Stratified 220-per-class sampling}
The canonical fit samples 220 pixels per admissible class without
replacement. This caps the per-class effective budget regardless of
how many labelled pixels a class actually has. Classes with very few
labelled pixels (alfalfa on Indian Pines, oats) are dropped from the
canonical fit on the standard $\ge 100$-pixel admissibility test.
The framework reports the admissibility decisions in the per-scene
JSON metadata.

\section{Empirical caveats specific to the benchmark set}

\subsection{Indian Pines is small and class-imbalanced}
At $145 \times 145$ with 16 classes and a long tail (some classes have
< 100 labelled pixels), Indian Pines is the smallest of the six
labelled scenes. The per-axis numbers on Indian Pines should be
treated with proportionate caution.

\subsection{Salinas-A is a subset of Salinas}
Salinas-A's six classes are a strict subset of Salinas's sixteen,
and the spatial footprint of Salinas-A is a contiguous rectangle
within Salinas. Numbers across these two ``scenes'' are
\emph{not} independent. A future paper may report only one of the
two; we report both because the band-mask diagnostic surfaces a
qualitatively different finding on Salinas-A (SWIR-only paired ARI
$= 0.766$) than on Salinas (SWIR-only paired ARI $= 0.435$), and
the contrast is itself a finding.

\subsection{Pavia U is VNIR-only by sensor design}
The ROSIS-03 sensor on Pavia U covers 430--860~nm, which is entirely
VNIR. The SWIR mask consequently retains 0 bands and is
auto-skipped. The no-water mask is a no-op on Pavia U (no bands
fall in the water-vapour ranges) and produces the same paired ARI as
the VNIR mask. These artefacts are recorded transparently in the
per-(scene, mask) JSON outputs.

\subsection{Botswana and KSC have wetland-dominant land cover}
The two scenes that consistently appear in the band-fragile regime
share a wetland-dominant land cover. Whether band-fragility is a
property of the LDA fit on these scenes, of the canonical
hyperparameter slate, or of the wetland-spectral-signature itself is
not resolvable from the present data; we flag the open question.

\section{Statistical caveats}

\subsection{Hierarchical Bayesian model assumes Gaussian residuals}
The F-1 likelihood
$\epsilon_{m,s,f} \sim \mathcal{N}(0, \sigma^2)$ is a standard
working assumption for OA-scaled residuals on $[0, 1]$. On the
present benchmarks the residuals are not visibly non-Gaussian (no
heavy tails, no bimodality observed in the posterior predictive)
but a Student-$t$ alternative may be appropriate on a benchmark set
with a strong outlier.

\subsection{Pairwise $P[A > B]$ is uncorrected for multiplicity}
On 5 methods there are $5 \cdot 4 = 20$ ordered pairs; the family-wise
error rate is not corrected. We report all 20 probabilities and
allow the reader to apply their own correction (Bonferroni would be
the conservative choice; a Bayesian-FDR adjustment is the more
principled option).

\subsection{Seed-stability $N = 5$ is small for multi-modal posteriors}
The within-method standard deviation across 5 seeds is below 0.03 on
all reported cells, suggesting a unimodal posterior on the present
benchmarks. KSC and Botswana, on which the canonical basis appears
weakly identified, may have multi-modal posteriors that the 5-seed
budget would not reveal. The companion code repository's
\texttt{deep\_seed\_stability/} block runs $N = 30$ seeds; the
framework's $N = 5$ choice is for parsimony with an explicit upgrade
path.

\subsection{$c_v$ co-occurrence statistics are computed on the same
corpus}
The $c_v$ reference corpus for the co-occurrence count is the
canonical document-term matrix of the same scene. Computing $c_v$ on
an external reference corpus (a published HSI library, for example)
would yield a different number and is a worthwhile follow-up.

\section{Reporting choices}

\subsection{Per-axis JSON over a single combined CSV}
The framework ships one JSON per axis per scene rather than a single
combined CSV. This is by design (clean separation of provenance) but
creates a slight friction for follow-up analysis that wants to
join across axes. The companion code repository provides a
\texttt{manifests/} index that lists every artefact path and its
schema, intended for join-on-load workflows.

\subsection{Figure choices in the main paper}
The journal main paper carries 4 figures (paired ARI heatmap,
Hungarian alignment, Bayesian forest plot, basis spectra grid). The
supplement carries the remaining 5 (coherence-vs-ARI, seed-stability
swarm, capacity sweep, cross-method 8x8 grid, HIDSAG band-mask). The
split reflects journal-section relevance, not figure quality
ordering.

\subsection{No statistical significance test for paired ARI}
Paired ARI does not carry a built-in significance threshold. We do
not compute permutation-test $p$-values per (scene, mask) tuple
because the test is computationally expensive on $D \sim 3000$
pixels and the framework's per-(scene, mask) values are already
extreme (paired ARI $\approx 0.01$ on KSC vs $0.77$ on Salinas-A
SWIR). A future paper applying the framework on a marginal-effect
benchmark would benefit from a permutation companion.

\section{What we would change in a 2nd-paper redesign}

If we were to redesign the framework from scratch given everything
we now know:

\begin{enumerate}
\item Move per-axis $K$-sweep from a separate axis (F-4) to a layer
that wraps every other axis: report (F-1, F-2, ..., F-12) at
$K \in \{4, 6, 8, 10, 12, 16\}$ rather than at the canonical $K$
only.
\item Replace F-12 (external baseline OA delta) with an
internal-comparison axis: compare the canonical topic-routed-soft
against a fixed deep baseline trained inside the framework on the
same train/test splits. The current external baseline carries the
cross-paper-comparison caveat that the deltas are not interpretable
without the per-paper split.
\item Add a per-pixel-confidence axis: report
$\mathrm{Pr}[\text{argmax}_k \theta_d = z_d^* \mid \theta_d]$
calibrated against a held-out validation fold. This would expose
the calibration of the topic representation directly.
\end{enumerate}

None of these changes invalidate the present framework; they
identify directions for forward-compatible extensions.

\bibliographystyle{IEEEtran}
\bibliography{../../bibliography/refs}

\end{document}
